\documentclass[12pt,a4paper]{article}
\usepackage[top=1.8cm,bottom=1.8cm,left=1.8cm,right=1.8cm]{geometry}
\usepackage{fontspec}
\setmainfont{TeX Gyre Termes}
\usepackage{amsmath,amssymb}
\usepackage{fancyhdr}
\usepackage{tasks}
\usepackage{tikz}
\settasks{label=\textbf{\Alph*.},label-width=1.5em,item-indent=2em}
\setlength{\parindent}{0pt}
\setlength{\parskip}{3pt}
\pagestyle{fancy}\fancyhf{}\renewcommand{\headrulewidth}{0pt}
\fancyhead[L]{\small Ngân hàng bài tập Toán 9 -- 9H1-C2-B4-D02}
\cfoot{\small Trang \thepage}
\newif\ifloigiai
\loigiaitrue
\newenvironment{cauhoi}[2]{%
  \par\medskip\noindent\rule{\textwidth}{0.4pt}\par\nopagebreak
  \noindent\textbf{[#1]}\hfill{\scriptsize\ttfamily #2}\par\nopagebreak\smallskip}{\par}
\newcommand{\loigiai}[1]{\ifloigiai\par\smallskip\noindent\textit{Lời giải.} #1\par\fi}
\newcommand{\dapan}[1]{\ifloigiai\par\smallskip\noindent\textit{Đáp án:} \textbf{#1}\par\fi}
\begin{document}
\begin{center}
{\Large\textbf{NGÂN HÀNG BÀI TẬP TOÁN 9}}\\[3pt]
\textbf{Chương 2. Phương trình và bất phương trình bậc nhất một ẩn}\\[2pt]
\textbf{Bài 4. Phương trình quy về phương trình bậc nhất một ẩn}\\[2pt]
\textbf{Dạng D02: Phương trình chứa ẩn ở mẫu}\\[2pt]
\textit{Mã dạng: 9H1-C2-B4-D02 \quad 50 bài tự luận + 50 câu trắc nghiệm \quad (TH: 25, VD: 25)}
\end{center}
\textbf{Quy ước mã câu:} \textit{khối, học kì, chương, bài, dạng, loại (TL/TN), mức độ, số thứ tự}.\par Ví dụ: \texttt{9H1-C2-B4-D02-TL-NB-001}.\par Để in bản cho học sinh, đặt \verb|\loigiaifalse| ngay sau \verb|\loigiaitrue|.
\section*{Phần I. Tự luận}
\begin{cauhoi}{Bài 1}{9H1-C2-B4-D02-TL-TH-001}
Giải phương trình $\dfrac{x+2}{x-4}=\dfrac{2}{5}$ ($x\ne 4$).
\loigiai{Điều kiện $x\ne 4$. Nhân hai vế với $x-4$: $x+2=\dfrac{2}{5}(x-4)$. Giải ta được $x=-6$ (thoả điều kiện).}
\end{cauhoi}
\begin{cauhoi}{Bài 2}{9H1-C2-B4-D02-TL-TH-002}
Giải phương trình $\dfrac{x+1}{x-1}=\dfrac{5}{3}$ ($x\ne 1$).
\loigiai{Điều kiện $x\ne 1$. Nhân hai vế với $x-1$: $x+1=\dfrac{5}{3}(x-1)$. Giải ta được $x=4$ (thoả điều kiện).}
\end{cauhoi}
\begin{cauhoi}{Bài 3}{9H1-C2-B4-D02-TL-TH-003}
Giải phương trình $\dfrac{x+1}{x-8}=-\dfrac{5}{4}$ ($x\ne 8$).
\loigiai{Điều kiện $x\ne 8$. Nhân hai vế với $x-8$: $x+1=-\dfrac{5}{4}(x-8)$. Giải ta được $x=4$ (thoả điều kiện).}
\end{cauhoi}
\begin{cauhoi}{Bài 4}{9H1-C2-B4-D02-TL-TH-004}
Giải phương trình $\dfrac{x+5}{x-2}=-6$ ($x\ne 2$).
\loigiai{Điều kiện $x\ne 2$. Nhân hai vế với $x-2$: $x+5=-6(x-2)$. Giải ta được $x=1$ (thoả điều kiện).}
\end{cauhoi}
\begin{cauhoi}{Bài 5}{9H1-C2-B4-D02-TL-TH-005}
Giải phương trình $\dfrac{x+4}{x-5}=0$ ($x\ne 5$).
\loigiai{Điều kiện $x\ne 5$. Nhân hai vế với $x-5$: $x+4=0(x-5)$. Giải ta được $x=-4$ (thoả điều kiện).}
\end{cauhoi}
\begin{cauhoi}{Bài 6}{9H1-C2-B4-D02-TL-TH-006}
Giải phương trình $\dfrac{x+3}{x-5}=-\dfrac{1}{3}$ ($x\ne 5$).
\loigiai{Điều kiện $x\ne 5$. Nhân hai vế với $x-5$: $x+3=-\dfrac{1}{3}(x-5)$. Giải ta được $x=-1$ (thoả điều kiện).}
\end{cauhoi}
\begin{cauhoi}{Bài 7}{9H1-C2-B4-D02-TL-TH-007}
Giải phương trình $\dfrac{x+5}{x-2}=-\dfrac{3}{4}$ ($x\ne 2$).
\loigiai{Điều kiện $x\ne 2$. Nhân hai vế với $x-2$: $x+5=-\dfrac{3}{4}(x-2)$. Giải ta được $x=-2$ (thoả điều kiện).}
\end{cauhoi}
\begin{cauhoi}{Bài 8}{9H1-C2-B4-D02-TL-TH-008}
Giải phương trình $\dfrac{x+4}{x-5}=10$ ($x\ne 5$).
\loigiai{Điều kiện $x\ne 5$. Nhân hai vế với $x-5$: $x+4=10(x-5)$. Giải ta được $x=6$ (thoả điều kiện).}
\end{cauhoi}
\begin{cauhoi}{Bài 9}{9H1-C2-B4-D02-TL-TH-009}
Giải phương trình $\dfrac{x+4}{x-2}=\dfrac{1}{7}$ ($x\ne 2$).
\loigiai{Điều kiện $x\ne 2$. Nhân hai vế với $x-2$: $x+4=\dfrac{1}{7}(x-2)$. Giải ta được $x=-5$ (thoả điều kiện).}
\end{cauhoi}
\begin{cauhoi}{Bài 10}{9H1-C2-B4-D02-TL-TH-010}
Giải phương trình $\dfrac{x+2}{x-9}=\dfrac{3}{14}$ ($x\ne 9$).
\loigiai{Điều kiện $x\ne 9$. Nhân hai vế với $x-9$: $x+2=\dfrac{3}{14}(x-9)$. Giải ta được $x=-5$ (thoả điều kiện).}
\end{cauhoi}
\begin{cauhoi}{Bài 11}{9H1-C2-B4-D02-TL-TH-011}
Giải phương trình $\dfrac{x+3}{x-7}=11$ ($x\ne 7$).
\loigiai{Điều kiện $x\ne 7$. Nhân hai vế với $x-7$: $x+3=11(x-7)$. Giải ta được $x=8$ (thoả điều kiện).}
\end{cauhoi}
\begin{cauhoi}{Bài 12}{9H1-C2-B4-D02-TL-TH-012}
Giải phương trình $\dfrac{x+3}{x-8}=\dfrac{2}{13}$ ($x\ne 8$).
\loigiai{Điều kiện $x\ne 8$. Nhân hai vế với $x-8$: $x+3=\dfrac{2}{13}(x-8)$. Giải ta được $x=-5$ (thoả điều kiện).}
\end{cauhoi}
\begin{cauhoi}{Bài 13}{9H1-C2-B4-D02-TL-TH-013}
Giải phương trình $\dfrac{x+5}{x-4}=-8$ ($x\ne 4$).
\loigiai{Điều kiện $x\ne 4$. Nhân hai vế với $x-4$: $x+5=-8(x-4)$. Giải ta được $x=3$ (thoả điều kiện).}
\end{cauhoi}
\begin{cauhoi}{Bài 14}{9H1-C2-B4-D02-TL-TH-014}
Giải phương trình $\dfrac{x+5}{x-7}=-3$ ($x\ne 7$).
\loigiai{Điều kiện $x\ne 7$. Nhân hai vế với $x-7$: $x+5=-3(x-7)$. Giải ta được $x=4$ (thoả điều kiện).}
\end{cauhoi}
\begin{cauhoi}{Bài 15}{9H1-C2-B4-D02-TL-TH-015}
Giải phương trình $\dfrac{x+5}{x-2}=\dfrac{1}{8}$ ($x\ne 2$).
\loigiai{Điều kiện $x\ne 2$. Nhân hai vế với $x-2$: $x+5=\dfrac{1}{8}(x-2)$. Giải ta được $x=-6$ (thoả điều kiện).}
\end{cauhoi}
\begin{cauhoi}{Bài 16}{9H1-C2-B4-D02-TL-TH-016}
Giải phương trình $\dfrac{x+4}{x-2}=\dfrac{1}{4}$ ($x\ne 2$).
\loigiai{Điều kiện $x\ne 2$. Nhân hai vế với $x-2$: $x+4=\dfrac{1}{4}(x-2)$. Giải ta được $x=-6$ (thoả điều kiện).}
\end{cauhoi}
\begin{cauhoi}{Bài 17}{9H1-C2-B4-D02-TL-TH-017}
Giải phương trình $\dfrac{x+3}{x-8}=-\dfrac{2}{9}$ ($x\ne 8$).
\loigiai{Điều kiện $x\ne 8$. Nhân hai vế với $x-8$: $x+3=-\dfrac{2}{9}(x-8)$. Giải ta được $x=-1$ (thoả điều kiện).}
\end{cauhoi}
\begin{cauhoi}{Bài 18}{9H1-C2-B4-D02-TL-TH-018}
Giải phương trình $\dfrac{x+8}{x-2}=\dfrac{7}{2}$ ($x\ne 2$).
\loigiai{Điều kiện $x\ne 2$. Nhân hai vế với $x-2$: $x+8=\dfrac{7}{2}(x-2)$. Giải ta được $x=6$ (thoả điều kiện).}
\end{cauhoi}
\begin{cauhoi}{Bài 19}{9H1-C2-B4-D02-TL-TH-019}
Giải phương trình $\dfrac{x+6}{x-2}=-7$ ($x\ne 2$).
\loigiai{Điều kiện $x\ne 2$. Nhân hai vế với $x-2$: $x+6=-7(x-2)$. Giải ta được $x=1$ (thoả điều kiện).}
\end{cauhoi}
\begin{cauhoi}{Bài 20}{9H1-C2-B4-D02-TL-TH-020}
Giải phương trình $\dfrac{x+6}{x-6}=13$ ($x\ne 6$).
\loigiai{Điều kiện $x\ne 6$. Nhân hai vế với $x-6$: $x+6=13(x-6)$. Giải ta được $x=7$ (thoả điều kiện).}
\end{cauhoi}
\begin{cauhoi}{Bài 21}{9H1-C2-B4-D02-TL-TH-021}
Giải phương trình $\dfrac{x+7}{x-5}=-11$ ($x\ne 5$).
\loigiai{Điều kiện $x\ne 5$. Nhân hai vế với $x-5$: $x+7=-11(x-5)$. Giải ta được $x=4$ (thoả điều kiện).}
\end{cauhoi}
\begin{cauhoi}{Bài 22}{9H1-C2-B4-D02-TL-TH-022}
Giải phương trình $\dfrac{x+8}{x-5}=\dfrac{16}{3}$ ($x\ne 5$).
\loigiai{Điều kiện $x\ne 5$. Nhân hai vế với $x-5$: $x+8=\dfrac{16}{3}(x-5)$. Giải ta được $x=8$ (thoả điều kiện).}
\end{cauhoi}
\begin{cauhoi}{Bài 23}{9H1-C2-B4-D02-TL-TH-023}
Giải phương trình $\dfrac{x+8}{x-4}=-\dfrac{7}{5}$ ($x\ne 4$).
\loigiai{Điều kiện $x\ne 4$. Nhân hai vế với $x-4$: $x+8=-\dfrac{7}{5}(x-4)$. Giải ta được $x=-1$ (thoả điều kiện).}
\end{cauhoi}
\begin{cauhoi}{Bài 24}{9H1-C2-B4-D02-TL-TH-024}
Giải phương trình $\dfrac{x+9}{x-9}=-\dfrac{1}{2}$ ($x\ne 9$).
\loigiai{Điều kiện $x\ne 9$. Nhân hai vế với $x-9$: $x+9=-\dfrac{1}{2}(x-9)$. Giải ta được $x=-3$ (thoả điều kiện).}
\end{cauhoi}
\begin{cauhoi}{Bài 25}{9H1-C2-B4-D02-TL-TH-025}
Giải phương trình $\dfrac{x+5}{x-3}=-7$ ($x\ne 3$).
\loigiai{Điều kiện $x\ne 3$. Nhân hai vế với $x-3$: $x+5=-7(x-3)$. Giải ta được $x=2$ (thoả điều kiện).}
\end{cauhoi}
\begin{cauhoi}{Bài 26}{9H1-C2-B4-D02-TL-VD-026}
Giải phương trình $\dfrac{x+3}{x-5}=9$ ($x\ne 5$).
\loigiai{Điều kiện $x\ne 5$. Nhân hai vế với $x-5$: $x+3=9(x-5)$. Giải ta được $x=6$ (thoả điều kiện).}
\end{cauhoi}
\begin{cauhoi}{Bài 27}{9H1-C2-B4-D02-TL-VD-027}
Giải phương trình $\dfrac{x+4}{x-1}=\dfrac{11}{6}$ ($x\ne 1$).
\loigiai{Điều kiện $x\ne 1$. Nhân hai vế với $x-1$: $x+4=\dfrac{11}{6}(x-1)$. Giải ta được $x=7$ (thoả điều kiện).}
\end{cauhoi}
\begin{cauhoi}{Bài 28}{9H1-C2-B4-D02-TL-VD-028}
Giải phương trình $\dfrac{x+8}{x-4}=-1$ ($x\ne 4$).
\loigiai{Điều kiện $x\ne 4$. Nhân hai vế với $x-4$: $x+8=-(x-4)$. Giải ta được $x=-2$ (thoả điều kiện).}
\end{cauhoi}
\begin{cauhoi}{Bài 29}{9H1-C2-B4-D02-TL-VD-029}
Giải phương trình $\dfrac{x+3}{x-6}=-\dfrac{1}{8}$ ($x\ne 6$).
\loigiai{Điều kiện $x\ne 6$. Nhân hai vế với $x-6$: $x+3=-\dfrac{1}{8}(x-6)$. Giải ta được $x=-2$ (thoả điều kiện).}
\end{cauhoi}
\begin{cauhoi}{Bài 30}{9H1-C2-B4-D02-TL-VD-030}
Giải phương trình $\dfrac{x+5}{x-1}=\dfrac{11}{5}$ ($x\ne 1$).
\loigiai{Điều kiện $x\ne 1$. Nhân hai vế với $x-1$: $x+5=\dfrac{11}{5}(x-1)$. Giải ta được $x=6$ (thoả điều kiện).}
\end{cauhoi}
\begin{cauhoi}{Bài 31}{9H1-C2-B4-D02-TL-VD-031}
Giải phương trình $\dfrac{x+9}{x-6}=-\dfrac{2}{3}$ ($x\ne 6$).
\loigiai{Điều kiện $x\ne 6$. Nhân hai vế với $x-6$: $x+9=-\dfrac{2}{3}(x-6)$. Giải ta được $x=-3$ (thoả điều kiện).}
\end{cauhoi}
\begin{cauhoi}{Bài 32}{9H1-C2-B4-D02-TL-VD-032}
Giải phương trình $\dfrac{x+2}{x-1}=\dfrac{8}{5}$ ($x\ne 1$).
\loigiai{Điều kiện $x\ne 1$. Nhân hai vế với $x-1$: $x+2=\dfrac{8}{5}(x-1)$. Giải ta được $x=6$ (thoả điều kiện).}
\end{cauhoi}
\begin{cauhoi}{Bài 33}{9H1-C2-B4-D02-TL-VD-033}
Giải phương trình $\dfrac{x+6}{x-3}=4$ ($x\ne 3$).
\loigiai{Điều kiện $x\ne 3$. Nhân hai vế với $x-3$: $x+6=4(x-3)$. Giải ta được $x=6$ (thoả điều kiện).}
\end{cauhoi}
\begin{cauhoi}{Bài 34}{9H1-C2-B4-D02-TL-VD-034}
Giải phương trình $\dfrac{x+9}{x-9}=-\dfrac{1}{5}$ ($x\ne 9$).
\loigiai{Điều kiện $x\ne 9$. Nhân hai vế với $x-9$: $x+9=-\dfrac{1}{5}(x-9)$. Giải ta được $x=-6$ (thoả điều kiện).}
\end{cauhoi}
\begin{cauhoi}{Bài 35}{9H1-C2-B4-D02-TL-VD-035}
Giải phương trình $\dfrac{x+7}{x-7}=-\dfrac{1}{13}$ ($x\ne 7$).
\loigiai{Điều kiện $x\ne 7$. Nhân hai vế với $x-7$: $x+7=-\dfrac{1}{13}(x-7)$. Giải ta được $x=-6$ (thoả điều kiện).}
\end{cauhoi}
\begin{cauhoi}{Bài 36}{9H1-C2-B4-D02-TL-VD-036}
Giải phương trình $\dfrac{x+3}{x-8}=-\dfrac{5}{6}$ ($x\ne 8$).
\loigiai{Điều kiện $x\ne 8$. Nhân hai vế với $x-8$: $x+3=-\dfrac{5}{6}(x-8)$. Giải ta được $x=2$ (thoả điều kiện).}
\end{cauhoi}
\begin{cauhoi}{Bài 37}{9H1-C2-B4-D02-TL-VD-037}
Giải phương trình $\dfrac{x+3}{x-8}=-\dfrac{4}{7}$ ($x\ne 8$).
\loigiai{Điều kiện $x\ne 8$. Nhân hai vế với $x-8$: $x+3=-\dfrac{4}{7}(x-8)$. Giải ta được $x=1$ (thoả điều kiện).}
\end{cauhoi}
\begin{cauhoi}{Bài 38}{9H1-C2-B4-D02-TL-VD-038}
Giải phương trình $\dfrac{x+5}{x-5}=\dfrac{13}{3}$ ($x\ne 5$).
\loigiai{Điều kiện $x\ne 5$. Nhân hai vế với $x-5$: $x+5=\dfrac{13}{3}(x-5)$. Giải ta được $x=8$ (thoả điều kiện).}
\end{cauhoi}
\begin{cauhoi}{Bài 39}{9H1-C2-B4-D02-TL-VD-039}
Giải phương trình $\dfrac{x+3}{x-7}=\dfrac{3}{13}$ ($x\ne 7$).
\loigiai{Điều kiện $x\ne 7$. Nhân hai vế với $x-7$: $x+3=\dfrac{3}{13}(x-7)$. Giải ta được $x=-6$ (thoả điều kiện).}
\end{cauhoi}
\begin{cauhoi}{Bài 40}{9H1-C2-B4-D02-TL-VD-040}
Giải phương trình $\dfrac{x+8}{x-4}=13$ ($x\ne 4$).
\loigiai{Điều kiện $x\ne 4$. Nhân hai vế với $x-4$: $x+8=13(x-4)$. Giải ta được $x=5$ (thoả điều kiện).}
\end{cauhoi}
\begin{cauhoi}{Bài 41}{9H1-C2-B4-D02-TL-VD-041}
Giải phương trình $\dfrac{x+3}{x-5}=-1$ ($x\ne 5$).
\loigiai{Điều kiện $x\ne 5$. Nhân hai vế với $x-5$: $x+3=-(x-5)$. Giải ta được $x=1$ (thoả điều kiện).}
\end{cauhoi}
\begin{cauhoi}{Bài 42}{9H1-C2-B4-D02-TL-VD-042}
Giải phương trình $\dfrac{x+1}{x-4}=6$ ($x\ne 4$).
\loigiai{Điều kiện $x\ne 4$. Nhân hai vế với $x-4$: $x+1=6(x-4)$. Giải ta được $x=5$ (thoả điều kiện).}
\end{cauhoi}
\begin{cauhoi}{Bài 43}{9H1-C2-B4-D02-TL-VD-043}
Giải phương trình $\dfrac{x+7}{x-1}=\dfrac{13}{5}$ ($x\ne 1$).
\loigiai{Điều kiện $x\ne 1$. Nhân hai vế với $x-1$: $x+7=\dfrac{13}{5}(x-1)$. Giải ta được $x=6$ (thoả điều kiện).}
\end{cauhoi}
\begin{cauhoi}{Bài 44}{9H1-C2-B4-D02-TL-VD-044}
Giải phương trình $\dfrac{x+4}{x-7}=\dfrac{3}{14}$ ($x\ne 7$).
\loigiai{Điều kiện $x\ne 7$. Nhân hai vế với $x-7$: $x+4=\dfrac{3}{14}(x-7)$. Giải ta được $x=-7$ (thoả điều kiện).}
\end{cauhoi}
\begin{cauhoi}{Bài 45}{9H1-C2-B4-D02-TL-VD-045}
Giải phương trình $\dfrac{x+4}{x-7}=-\dfrac{1}{10}$ ($x\ne 7$).
\loigiai{Điều kiện $x\ne 7$. Nhân hai vế với $x-7$: $x+4=-\dfrac{1}{10}(x-7)$. Giải ta được $x=-3$ (thoả điều kiện).}
\end{cauhoi}
\begin{cauhoi}{Bài 46}{9H1-C2-B4-D02-TL-VD-046}
Giải phương trình $\dfrac{x+2}{x-7}=-\dfrac{4}{5}$ ($x\ne 7$).
\loigiai{Điều kiện $x\ne 7$. Nhân hai vế với $x-7$: $x+2=-\dfrac{4}{5}(x-7)$. Giải ta được $x=2$ (thoả điều kiện).}
\end{cauhoi}
\begin{cauhoi}{Bài 47}{9H1-C2-B4-D02-TL-VD-047}
Giải phương trình $\dfrac{x+9}{x-2}=-\dfrac{2}{9}$ ($x\ne 2$).
\loigiai{Điều kiện $x\ne 2$. Nhân hai vế với $x-2$: $x+9=-\dfrac{2}{9}(x-2)$. Giải ta được $x=-7$ (thoả điều kiện).}
\end{cauhoi}
\begin{cauhoi}{Bài 48}{9H1-C2-B4-D02-TL-VD-048}
Giải phương trình $\dfrac{x+6}{x-8}=\dfrac{1}{8}$ ($x\ne 8$).
\loigiai{Điều kiện $x\ne 8$. Nhân hai vế với $x-8$: $x+6=\dfrac{1}{8}(x-8)$. Giải ta được $x=-8$ (thoả điều kiện).}
\end{cauhoi}
\begin{cauhoi}{Bài 49}{9H1-C2-B4-D02-TL-VD-049}
Giải phương trình $\dfrac{x+9}{x-1}=\dfrac{7}{2}$ ($x\ne 1$).
\loigiai{Điều kiện $x\ne 1$. Nhân hai vế với $x-1$: $x+9=\dfrac{7}{2}(x-1)$. Giải ta được $x=5$ (thoả điều kiện).}
\end{cauhoi}
\begin{cauhoi}{Bài 50}{9H1-C2-B4-D02-TL-VD-050}
Giải phương trình $\dfrac{x+4}{x-3}=-\dfrac{3}{4}$ ($x\ne 3$).
\loigiai{Điều kiện $x\ne 3$. Nhân hai vế với $x-3$: $x+4=-\dfrac{3}{4}(x-3)$. Giải ta được $x=-1$ (thoả điều kiện).}
\end{cauhoi}
\clearpage\section*{Phần II. Trắc nghiệm}
\begin{cauhoi}{Câu 1}{9H1-C2-B4-D02-TN-TH-001}
Giải phương trình $\dfrac{x+2}{x-4}=\dfrac{2}{5}$ ($x\ne 4$).
\begin{tasks}(2)
\task $-6$
\task $6$
\task $4$
\task $-5$
\end{tasks}
\dapan{A}
\end{cauhoi}
\begin{cauhoi}{Câu 2}{9H1-C2-B4-D02-TN-TH-002}
Giải phương trình $\dfrac{x+1}{x-1}=\dfrac{5}{3}$ ($x\ne 1$).
\begin{tasks}(2)
\task $4$
\task $5$
\task $1$
\task $-4$
\end{tasks}
\dapan{A}
\end{cauhoi}
\begin{cauhoi}{Câu 3}{9H1-C2-B4-D02-TN-TH-003}
Giải phương trình $\dfrac{x+1}{x-8}=-\dfrac{5}{4}$ ($x\ne 8$).
\begin{tasks}(2)
\task $8$
\task $5$
\task $-4$
\task $4$
\end{tasks}
\dapan{D}
\end{cauhoi}
\begin{cauhoi}{Câu 4}{9H1-C2-B4-D02-TN-TH-004}
Giải phương trình $\dfrac{x+5}{x-2}=-6$ ($x\ne 2$).
\begin{tasks}(2)
\task $0$
\task $2$
\task $-1$
\task $1$
\end{tasks}
\dapan{D}
\end{cauhoi}
\begin{cauhoi}{Câu 5}{9H1-C2-B4-D02-TN-TH-005}
Giải phương trình $\dfrac{x+4}{x-5}=0$ ($x\ne 5$).
\begin{tasks}(2)
\task $-3$
\task $4$
\task $-4$
\task $5$
\end{tasks}
\dapan{C}
\end{cauhoi}
\begin{cauhoi}{Câu 6}{9H1-C2-B4-D02-TN-TH-006}
Giải phương trình $\dfrac{x+3}{x-5}=-\dfrac{1}{3}$ ($x\ne 5$).
\begin{tasks}(2)
\task $1$
\task $-1$
\task $5$
\task $0$
\end{tasks}
\dapan{B}
\end{cauhoi}
\begin{cauhoi}{Câu 7}{9H1-C2-B4-D02-TN-TH-007}
Giải phương trình $\dfrac{x+5}{x-2}=-\dfrac{3}{4}$ ($x\ne 2$).
\begin{tasks}(2)
\task $-1$
\task $2$
\task $-3$
\task $-2$
\end{tasks}
\dapan{D}
\end{cauhoi}
\begin{cauhoi}{Câu 8}{9H1-C2-B4-D02-TN-TH-008}
Giải phương trình $\dfrac{x+4}{x-5}=10$ ($x\ne 5$).
\begin{tasks}(2)
\task $5$
\task $7$
\task $6$
\task $-6$
\end{tasks}
\dapan{C}
\end{cauhoi}
\begin{cauhoi}{Câu 9}{9H1-C2-B4-D02-TN-TH-009}
Giải phương trình $\dfrac{x+4}{x-2}=\dfrac{1}{7}$ ($x\ne 2$).
\begin{tasks}(2)
\task $5$
\task $-5$
\task $2$
\task $-4$
\end{tasks}
\dapan{B}
\end{cauhoi}
\begin{cauhoi}{Câu 10}{9H1-C2-B4-D02-TN-TH-010}
Giải phương trình $\dfrac{x+2}{x-9}=\dfrac{3}{14}$ ($x\ne 9$).
\begin{tasks}(2)
\task $-5$
\task $9$
\task $-4$
\task $5$
\end{tasks}
\dapan{A}
\end{cauhoi}
\begin{cauhoi}{Câu 11}{9H1-C2-B4-D02-TN-TH-011}
Giải phương trình $\dfrac{x+3}{x-7}=11$ ($x\ne 7$).
\begin{tasks}(2)
\task $8$
\task $-8$
\task $9$
\task $7$
\end{tasks}
\dapan{A}
\end{cauhoi}
\begin{cauhoi}{Câu 12}{9H1-C2-B4-D02-TN-TH-012}
Giải phương trình $\dfrac{x+3}{x-8}=\dfrac{2}{13}$ ($x\ne 8$).
\begin{tasks}(2)
\task $8$
\task $-5$
\task $-4$
\task $5$
\end{tasks}
\dapan{B}
\end{cauhoi}
\begin{cauhoi}{Câu 13}{9H1-C2-B4-D02-TN-TH-013}
Giải phương trình $\dfrac{x+5}{x-4}=-8$ ($x\ne 4$).
\begin{tasks}(2)
\task $4$
\task $2$
\task $3$
\task $-3$
\end{tasks}
\dapan{C}
\end{cauhoi}
\begin{cauhoi}{Câu 14}{9H1-C2-B4-D02-TN-TH-014}
Giải phương trình $\dfrac{x+5}{x-7}=-3$ ($x\ne 7$).
\begin{tasks}(2)
\task $-4$
\task $7$
\task $4$
\task $5$
\end{tasks}
\dapan{C}
\end{cauhoi}
\begin{cauhoi}{Câu 15}{9H1-C2-B4-D02-TN-TH-015}
Giải phương trình $\dfrac{x+5}{x-2}=\dfrac{1}{8}$ ($x\ne 2$).
\begin{tasks}(2)
\task $6$
\task $-6$
\task $-5$
\task $2$
\end{tasks}
\dapan{B}
\end{cauhoi}
\begin{cauhoi}{Câu 16}{9H1-C2-B4-D02-TN-TH-016}
Giải phương trình $\dfrac{x+4}{x-2}=\dfrac{1}{4}$ ($x\ne 2$).
\begin{tasks}(2)
\task $-5$
\task $-6$
\task $6$
\task $2$
\end{tasks}
\dapan{B}
\end{cauhoi}
\begin{cauhoi}{Câu 17}{9H1-C2-B4-D02-TN-TH-017}
Giải phương trình $\dfrac{x+3}{x-8}=-\dfrac{2}{9}$ ($x\ne 8$).
\begin{tasks}(2)
\task $8$
\task $-1$
\task $0$
\task $1$
\end{tasks}
\dapan{B}
\end{cauhoi}
\begin{cauhoi}{Câu 18}{9H1-C2-B4-D02-TN-TH-018}
Giải phương trình $\dfrac{x+8}{x-2}=\dfrac{7}{2}$ ($x\ne 2$).
\begin{tasks}(2)
\task $-6$
\task $2$
\task $6$
\task $7$
\end{tasks}
\dapan{C}
\end{cauhoi}
\begin{cauhoi}{Câu 19}{9H1-C2-B4-D02-TN-TH-019}
Giải phương trình $\dfrac{x+6}{x-2}=-7$ ($x\ne 2$).
\begin{tasks}(2)
\task $-1$
\task $1$
\task $0$
\task $2$
\end{tasks}
\dapan{B}
\end{cauhoi}
\begin{cauhoi}{Câu 20}{9H1-C2-B4-D02-TN-TH-020}
Giải phương trình $\dfrac{x+6}{x-6}=13$ ($x\ne 6$).
\begin{tasks}(2)
\task $8$
\task $-7$
\task $7$
\task $6$
\end{tasks}
\dapan{C}
\end{cauhoi}
\begin{cauhoi}{Câu 21}{9H1-C2-B4-D02-TN-TH-021}
Giải phương trình $\dfrac{x+7}{x-5}=-11$ ($x\ne 5$).
\begin{tasks}(2)
\task $5$
\task $3$
\task $-4$
\task $4$
\end{tasks}
\dapan{D}
\end{cauhoi}
\begin{cauhoi}{Câu 22}{9H1-C2-B4-D02-TN-TH-022}
Giải phương trình $\dfrac{x+8}{x-5}=\dfrac{16}{3}$ ($x\ne 5$).
\begin{tasks}(2)
\task $5$
\task $9$
\task $8$
\task $-8$
\end{tasks}
\dapan{C}
\end{cauhoi}
\begin{cauhoi}{Câu 23}{9H1-C2-B4-D02-TN-TH-023}
Giải phương trình $\dfrac{x+8}{x-4}=-\dfrac{7}{5}$ ($x\ne 4$).
\begin{tasks}(2)
\task $0$
\task $1$
\task $4$
\task $-1$
\end{tasks}
\dapan{D}
\end{cauhoi}
\begin{cauhoi}{Câu 24}{9H1-C2-B4-D02-TN-TH-024}
Giải phương trình $\dfrac{x+9}{x-9}=-\dfrac{1}{2}$ ($x\ne 9$).
\begin{tasks}(2)
\task $9$
\task $-3$
\task $-2$
\task $3$
\end{tasks}
\dapan{B}
\end{cauhoi}
\begin{cauhoi}{Câu 25}{9H1-C2-B4-D02-TN-TH-025}
Giải phương trình $\dfrac{x+5}{x-3}=-7$ ($x\ne 3$).
\begin{tasks}(2)
\task $2$
\task $3$
\task $-2$
\task $1$
\end{tasks}
\dapan{A}
\end{cauhoi}
\begin{cauhoi}{Câu 26}{9H1-C2-B4-D02-TN-VD-026}
Giải phương trình $\dfrac{x+3}{x-5}=9$ ($x\ne 5$).
\begin{tasks}(2)
\task $6$
\task $5$
\task $-6$
\task $7$
\end{tasks}
\dapan{A}
\end{cauhoi}
\begin{cauhoi}{Câu 27}{9H1-C2-B4-D02-TN-VD-027}
Giải phương trình $\dfrac{x+4}{x-1}=\dfrac{11}{6}$ ($x\ne 1$).
\begin{tasks}(2)
\task $-7$
\task $8$
\task $1$
\task $7$
\end{tasks}
\dapan{D}
\end{cauhoi}
\begin{cauhoi}{Câu 28}{9H1-C2-B4-D02-TN-VD-028}
Giải phương trình $\dfrac{x+8}{x-4}=-1$ ($x\ne 4$).
\begin{tasks}(2)
\task $4$
\task $2$
\task $-2$
\task $-1$
\end{tasks}
\dapan{C}
\end{cauhoi}
\begin{cauhoi}{Câu 29}{9H1-C2-B4-D02-TN-VD-029}
Giải phương trình $\dfrac{x+3}{x-6}=-\dfrac{1}{8}$ ($x\ne 6$).
\begin{tasks}(2)
\task $-2$
\task $6$
\task $-1$
\task $2$
\end{tasks}
\dapan{A}
\end{cauhoi}
\begin{cauhoi}{Câu 30}{9H1-C2-B4-D02-TN-VD-030}
Giải phương trình $\dfrac{x+5}{x-1}=\dfrac{11}{5}$ ($x\ne 1$).
\begin{tasks}(2)
\task $7$
\task $-6$
\task $1$
\task $6$
\end{tasks}
\dapan{D}
\end{cauhoi}
\begin{cauhoi}{Câu 31}{9H1-C2-B4-D02-TN-VD-031}
Giải phương trình $\dfrac{x+9}{x-6}=-\dfrac{2}{3}$ ($x\ne 6$).
\begin{tasks}(2)
\task $6$
\task $3$
\task $-2$
\task $-3$
\end{tasks}
\dapan{D}
\end{cauhoi}
\begin{cauhoi}{Câu 32}{9H1-C2-B4-D02-TN-VD-032}
Giải phương trình $\dfrac{x+2}{x-1}=\dfrac{8}{5}$ ($x\ne 1$).
\begin{tasks}(2)
\task $6$
\task $1$
\task $7$
\task $-6$
\end{tasks}
\dapan{A}
\end{cauhoi}
\begin{cauhoi}{Câu 33}{9H1-C2-B4-D02-TN-VD-033}
Giải phương trình $\dfrac{x+6}{x-3}=4$ ($x\ne 3$).
\begin{tasks}(2)
\task $3$
\task $7$
\task $6$
\task $-6$
\end{tasks}
\dapan{C}
\end{cauhoi}
\begin{cauhoi}{Câu 34}{9H1-C2-B4-D02-TN-VD-034}
Giải phương trình $\dfrac{x+9}{x-9}=-\dfrac{1}{5}$ ($x\ne 9$).
\begin{tasks}(2)
\task $-5$
\task $-6$
\task $9$
\task $6$
\end{tasks}
\dapan{B}
\end{cauhoi}
\begin{cauhoi}{Câu 35}{9H1-C2-B4-D02-TN-VD-035}
Giải phương trình $\dfrac{x+7}{x-7}=-\dfrac{1}{13}$ ($x\ne 7$).
\begin{tasks}(2)
\task $6$
\task $7$
\task $-5$
\task $-6$
\end{tasks}
\dapan{D}
\end{cauhoi}
\begin{cauhoi}{Câu 36}{9H1-C2-B4-D02-TN-VD-036}
Giải phương trình $\dfrac{x+3}{x-8}=-\dfrac{5}{6}$ ($x\ne 8$).
\begin{tasks}(2)
\task $2$
\task $3$
\task $8$
\task $-2$
\end{tasks}
\dapan{A}
\end{cauhoi}
\begin{cauhoi}{Câu 37}{9H1-C2-B4-D02-TN-VD-037}
Giải phương trình $\dfrac{x+3}{x-8}=-\dfrac{4}{7}$ ($x\ne 8$).
\begin{tasks}(2)
\task $1$
\task $8$
\task $2$
\task $-1$
\end{tasks}
\dapan{A}
\end{cauhoi}
\begin{cauhoi}{Câu 38}{9H1-C2-B4-D02-TN-VD-038}
Giải phương trình $\dfrac{x+5}{x-5}=\dfrac{13}{3}$ ($x\ne 5$).
\begin{tasks}(2)
\task $8$
\task $9$
\task $-8$
\task $5$
\end{tasks}
\dapan{A}
\end{cauhoi}
\begin{cauhoi}{Câu 39}{9H1-C2-B4-D02-TN-VD-039}
Giải phương trình $\dfrac{x+3}{x-7}=\dfrac{3}{13}$ ($x\ne 7$).
\begin{tasks}(2)
\task $7$
\task $-5$
\task $-6$
\task $6$
\end{tasks}
\dapan{C}
\end{cauhoi}
\begin{cauhoi}{Câu 40}{9H1-C2-B4-D02-TN-VD-040}
Giải phương trình $\dfrac{x+8}{x-4}=13$ ($x\ne 4$).
\begin{tasks}(2)
\task $-5$
\task $6$
\task $4$
\task $5$
\end{tasks}
\dapan{D}
\end{cauhoi}
\begin{cauhoi}{Câu 41}{9H1-C2-B4-D02-TN-VD-041}
Giải phương trình $\dfrac{x+3}{x-5}=-1$ ($x\ne 5$).
\begin{tasks}(2)
\task $5$
\task $1$
\task $2$
\task $-1$
\end{tasks}
\dapan{B}
\end{cauhoi}
\begin{cauhoi}{Câu 42}{9H1-C2-B4-D02-TN-VD-042}
Giải phương trình $\dfrac{x+1}{x-4}=6$ ($x\ne 4$).
\begin{tasks}(2)
\task $6$
\task $5$
\task $4$
\task $-5$
\end{tasks}
\dapan{B}
\end{cauhoi}
\begin{cauhoi}{Câu 43}{9H1-C2-B4-D02-TN-VD-043}
Giải phương trình $\dfrac{x+7}{x-1}=\dfrac{13}{5}$ ($x\ne 1$).
\begin{tasks}(2)
\task $-6$
\task $6$
\task $7$
\task $1$
\end{tasks}
\dapan{B}
\end{cauhoi}
\begin{cauhoi}{Câu 44}{9H1-C2-B4-D02-TN-VD-044}
Giải phương trình $\dfrac{x+4}{x-7}=\dfrac{3}{14}$ ($x\ne 7$).
\begin{tasks}(2)
\task $-6$
\task $-7$
\task $-8$
\task $7$
\end{tasks}
\dapan{B}
\end{cauhoi}
\begin{cauhoi}{Câu 45}{9H1-C2-B4-D02-TN-VD-045}
Giải phương trình $\dfrac{x+4}{x-7}=-\dfrac{1}{10}$ ($x\ne 7$).
\begin{tasks}(2)
\task $3$
\task $-2$
\task $-3$
\task $7$
\end{tasks}
\dapan{C}
\end{cauhoi}
\begin{cauhoi}{Câu 46}{9H1-C2-B4-D02-TN-VD-046}
Giải phương trình $\dfrac{x+2}{x-7}=-\dfrac{4}{5}$ ($x\ne 7$).
\begin{tasks}(2)
\task $-2$
\task $3$
\task $7$
\task $2$
\end{tasks}
\dapan{D}
\end{cauhoi}
\begin{cauhoi}{Câu 47}{9H1-C2-B4-D02-TN-VD-047}
Giải phương trình $\dfrac{x+9}{x-2}=-\dfrac{2}{9}$ ($x\ne 2$).
\begin{tasks}(2)
\task $7$
\task $-7$
\task $2$
\task $-6$
\end{tasks}
\dapan{B}
\end{cauhoi}
\begin{cauhoi}{Câu 48}{9H1-C2-B4-D02-TN-VD-048}
Giải phương trình $\dfrac{x+6}{x-8}=\dfrac{1}{8}$ ($x\ne 8$).
\begin{tasks}(2)
\task $8$
\task $-8$
\task $-9$
\task $-7$
\end{tasks}
\dapan{B}
\end{cauhoi}
\begin{cauhoi}{Câu 49}{9H1-C2-B4-D02-TN-VD-049}
Giải phương trình $\dfrac{x+9}{x-1}=\dfrac{7}{2}$ ($x\ne 1$).
\begin{tasks}(2)
\task $-5$
\task $5$
\task $6$
\task $1$
\end{tasks}
\dapan{B}
\end{cauhoi}
\begin{cauhoi}{Câu 50}{9H1-C2-B4-D02-TN-VD-050}
Giải phương trình $\dfrac{x+4}{x-3}=-\dfrac{3}{4}$ ($x\ne 3$).
\begin{tasks}(2)
\task $0$
\task $1$
\task $3$
\task $-1$
\end{tasks}
\dapan{D}
\end{cauhoi}
\clearpage\section*{Bảng đáp án trắc nghiệm}
\begin{center}\begin{tabular}{|c|c|c|c|c|c|c|c|c|c|}
\hline
1. A & 2. A & 3. D & 4. D & 5. C & 6. B & 7. D & 8. C & 9. B & 10. A\\\hline
11. A & 12. B & 13. C & 14. C & 15. B & 16. B & 17. B & 18. C & 19. B & 20. C\\\hline
21. D & 22. C & 23. D & 24. B & 25. A & 26. A & 27. D & 28. C & 29. A & 30. D\\\hline
31. D & 32. A & 33. C & 34. B & 35. D & 36. A & 37. A & 38. A & 39. C & 40. D\\\hline
41. B & 42. B & 43. B & 44. B & 45. C & 46. D & 47. B & 48. B & 49. B & 50. D\\\hline
\end{tabular}\end{center}
\end{document}
